\subsection{\texorpdfstring{空间 \(\mathcal{F}_{\mathrm{F}}^{p}\)}{空间 F sub F to the p}}

设 F 是一个 Banach 空间，\(\mathcal{F}(X;F)\) （或简作 \(F_{F}\) ）是 X 到 F 中所有映射所成的向量空间。对 \(1 \leqslant p < +\infty\) ，我们将用 \(\mathcal{F}^{p}(X,\mu;F)\) 或 \(\mathcal{F}_{\mathrm{F}}^{p}(X,\mu)\) ，或简作 \(\mathcal{F}_{\mathrm{F}}^{p}(\mu)\) ，或 \(\mathcal{F}_{\mathrm{F}}^{p}\) （若不致引起混淆），表示 X 到 F 中使 \(N_{p}(\mathbf{f}) < +\infty\) 的映射 f 所成之集（我们用 \(F^{p}\) 代替 \(F_{R}^{p}\) ）。显然 \(\mathcal{F}_{\mathrm{F}}^{p}(|\mu|) = \mathcal{F}_{\mathrm{F}}^{p}(\mu)\) 。由 No.~2 的命题 2 立得：\(F_{F}^{p}\) 是 \(F_{F}\) 的一个线性子空间，且 \(N_{p}(\mathbf{f})\) 是这个空间上的一个半范数。我们总假定（除明确声明相反情形外）\(F_{F}^{p}\) 赋有由这个半范数定义的拓扑；我们称这拓扑为 p 阶平均收敛拓扑（对 p = 1，简称平均收敛拓扑；对 p = 2，也称«均方收敛拓扑»）。我们称 \(F_{F}^{p}\) 上关于这拓扑收敛于 f 的滤子 G （相应地，序列 \((\mathbf{f}_{n})\) ，其元素属于 \(F_{F}^{p}\) ）为 p 阶平均收敛于 f 的；这就是说，\(N_{p}(\mathbf{g}-\mathbf{f})\) 关于 G 趋于 0 （相应地， \(N_{p}(\mathbf{f}_{n}-\mathbf{f})\) 当 n 趋于无穷时趋于 0）。

这一术语可直接推广到下述情形：函数 \(f_{n}\) 与函数 f 只是几乎处处有定义（或取值于 \(\overline{R}\) ，且几乎处处有定义并有限）。

注意，局部凸空间 \(F_{F}^{p}\) 一般不是 Hausdorff 的；0 在这个空间中的闭包是 X 到 F 中的可忽略映射所成的线性子空间 \(N_{F}\) （No.~2, 命题 3）。

注. --- 设 F 是复数域 C 上的一个 Banach 空间；于是，对每个函数 \(f \in F_{F}^{p}\) 与每个复数 \(\alpha\) ，\(\alpha f\) 属于 \(F_{F}^{p}\) 且 \(N_{p}(\alpha f) = |\alpha| N_{p}(f)\) ；换句话说，\(F_{F}^{p}\) 也是 C 上的一个向量空间，且 \(N_{p}(f)\) 是这个复向量空间上的一个半范数（TVS, II, §1）。

命题 4. --- 设 B 是 \(F_{F}^{p}\) 上的一个滤子基。假定存在紧集 \(K \subset X\) ，使对每个集 \(M \in B\) ，所有映射 \(f \in M\) 的支集都含于 K。在这些条件下，若 B 在 X 上一致收敛于 \(f_{0}\) ，则 \(f_{0}\) 属于 \(F_{F}^{p}\) 且 B p 阶平均收敛于 \(f_{0}\) 。

这等价于说：在支集含于一个固定紧集的映射 \(f \in F_{F}^{p}\) 所成之集上，一致收敛拓扑细于 p 阶平均收敛拓扑。

因为，设 \(h\) 是 \(X\) 到 \([0,1]\) 中的一个连续映射，具有紧支集，且在 \(K\) 上等于 1（Ch. III, §1, No.~2, 引理 1）。对每个 \(\varepsilon > 0\) ，存在 \(M \in \mathfrak{B}\) ，使对每个映射 \(\mathbf{f} \in M\) ，有 \(|\mathbf{f}(x) - \mathbf{f}_0(x)| \leqslant \varepsilon h(x)\) 对所有 \(x \in X\) 成立。由此得出 \(N sub p(\mathbf{f} - \mathbf{f}_0) \leqslant \varepsilon N sub p(h)\) ，命题得证。

命题 5. --- 局部凸空间 \(\mathcal{F}_{\mathrm{F}}^{p}\) 是完备的。

由于与 \(\mathcal{F}_{\mathrm{F}}^{p}\) 相伴的 Hausdorff 空间是一个赋范空间，只需证明每个 Cauchy 序列 \((\mathbf{f}_n)\) （\(\mathcal{F}_{\mathrm{F}}^{p}\) 中的）对于 \(p\) 阶平均收敛拓扑都有一个极限（GT, IX, §2, No.~6, 命题 9）。按假设，对每个 \(\varepsilon > 0\) 存在整数 \(m_0\) ，使关系 \(m \geqslant m_0\) 、 \(n \geqslant m_0\) 蕴涵 \(\mathrm{N}_p(\mathbf{f}_n - \mathbf{f}_m) \leqslant \varepsilon\) 。于是可以对 \(k\) 用归纳法定义一个严格递增的整数序列 \((n_k)\) （ \(\geqslant 0\) ），使 \(\mathrm{N}_p(\mathbf{f}_{n_{k+1}} - \mathbf{f}_{n_k}) \leqslant 2^{-k}\) 。若我们能证明通项为 \(\mathbf{g}_k = \mathbf{f}_{n_{k+1}} - \mathbf{f}_{n_k}\) （ \(k \geqslant 1\) ）的级数 \(p\) 阶平均收敛，则它将有和 \(\mathbf{g} \in \mathcal{F}_{\mathrm{F}}^p\) ，而 \(\mathbf{f} = \mathbf{g} + \mathbf{f}_{n_1}\) 将是序列 \((\mathbf{f}_{n_k})\) 在 \(\mathcal{F}_{\mathrm{F}}^p\) 中的极限；于是 \(\mathbf{f}\) 是序列 \((\mathbf{f}_n)\) 的一个接触点；由于这个序列是 Cauchy 序列，它将以 f 为极限，从而命题 5 得证（GT, II, §3, No.~2, 命题 5 的推论 2）。

于是命题 5 是下述命题的推论：

命题 6. --- 设 \((\mathbf{f}_{n})\) 是 \(F_{F}^{p}\) 中函数的一个序列，使 \(\sum_{n=1}^{\infty}\mathrm{N}_{p}(\mathbf{f}_{n})<+\infty\) 。在这些条件下，通项为 \(\mathbf{f}_{n}(x)\in\mathbf{F}\) 的级数在 X 上几乎处处绝对收敛。在级数收敛的点处令 \(\mathbf{f}(x)=\sum_{n=1}^{\infty}\mathbf{f}_{n}(x)\) ，在别处令 \(\mathbf{f}(x)=0\) ，则函数 f 属于 \(F_{F}^{p}\) 且是通项为 \(f_{n}\) 的级数的和（对于 p 阶平均收敛拓扑）；更精确地说，对每个 \(n\geqslant0\) ，

\[
\mathrm{N} _ {p} \left(\mathbf {f} - \sum_ {k = 1} ^ {n} \mathbf {f} _ {k}\right) \leqslant \sum_ {k = n + 1} ^ {\infty} \mathrm{N} _ {p} (\mathbf {f} _ {k}).\tag{8}
\]

考虑正函数（有限或非有限） \(g(x) = \sum_{n=1}^{\infty} |\mathbf{f}_n(x)|\) 。由可数凸性定理（No.~2, 定理 1），

\[
\mathrm{N} _ {p} (g) \leqslant \sum_ {n = 1} ^ {\infty} \mathrm{N} _ {p} (\mathbf {f} _ {n}) <   + \infty ;
\]

于是 \(g\) 几乎处处有限（§2, No.~3, 命题 7），这表明通项为 \(\mathbf{f}_n(x)\) 的级数几乎处处绝对收敛。由于 \(F\) 是完备的，这级数几乎处处收敛，且对每个 \(x \in X\) ，有 \(|\mathbf{f}(x)| \leqslant \sum_{n=1}^{\infty} |\mathbf{f}_n(x)| = g(x)\) ，由此得

\[
\mathrm{N} _ {p} (\mathbf {f}) \leqslant \mathrm{N} _ {p} (g) \leqslant \sum_ {n = 1} ^ {\infty} \mathrm{N} _ {p} (\mathbf {f} _ {n}) <   + \infty ,
\]

这就证明了 \(\mathbf{f}\) 属于 \(\mathcal{F}_{\mathrm{F}}^{p}\) 。另一方面，对每个整数 \(n\) ，

\[
\left| \mathbf {f} (x) - \sum_ {k = 1} ^ {n} \mathbf {f} _ {k} (x) \right| \leqslant \sum_ {k = n + 1} ^ {\infty} | \mathbf {f} _ {k} (x) |
\]

几乎处处成立，由此得 \(\mathrm{N}_{p}\left(\mathbf{f}-\sum_{k=1}^{n}\mathbf{f}_{k}\right)\leqslant\sum_{k=n+1}^{\infty}\mathrm{N}_{p}(\mathbf{f}_{k})\) 。按假设，通项为 \(\mathrm{N}_{p}(\mathbf{f}_{n})\) 的级数收敛；因此，对每个 \(\varepsilon > 0\) ，存在整数 \(n\) 使 \(\sum_{k=n+1}^{\infty} \mathrm{N}_{p}(\mathbf{f}_{k}) \leqslant \varepsilon\) ，而不等式 (8) 证明了 \(\mathbf{f}\) 是通项为 \(\mathbf{f}_{n}\) 的级数的和（对于 \(p\) 阶平均收敛拓扑）。

于是命题 5 与命题 6 完全得证。

